% -----------------------------------------------------------
% บทที่ 5: สรุปผลและข้อเสนอแนะ
% -----------------------------------------------------------

% -----------------------------------------------------------
\section{สรุปผล}
% -----------------------------------------------------------

\subsection{สรุปภาพรวมของโครงงาน}

โครงงานนี้เริ่มจากคำถามว่า ถ้านำโมเดลเปลี่ยนทรงผมที่มีอยู่แล้วมาปรับจูนด้วยข้อมูลของผู้วิจัยเอง
จะแก้จุดอ่อนที่เห็นได้หรือไม่ และโมเดลแบบใดเหมาะกับการใช้งานจริง
ผู้วิจัยเลือกโมเดลโครงข่ายปรปักษ์เชิงกำเนิด (Generative Adversarial Network: GAN) 2 ตัวที่ทำงานต่างกันอย่างชัดเจน
คือ HairFastGAN ที่ใช้ตัวเข้ารหัส (Encoder-based) และ Barbershop ที่ค้นหารหัสทีละภาพ (Optimization-based)
แล้วปรับจูนทั้งคู่ด้วยชุดฝึก 150 ชุด ใช้ฟังก์ชันสูญเสียหลักชุดเดียวกันซึ่งวัด 5 ด้าน
และทดลองอัตราการเรียนรู้ (Learning Rate) 4 ชุด (Trial) ชุดละ 30 Epoch
การเลือก Checkpoint ใช้กติกาฉบับแก้ไขที่เปิดเผยไว้ในบทที่ 3 คือคะแนนรวมร่วมกับเกณฑ์ 5 ด่านบนชุดตรวจสอบ
จากนั้นวัดผลบนชุดทดสอบ 10 ชุด เพิ่มขั้นขัดเกลาด้วย Stable Diffusion Inpainting แบบแรงต่ำ
ให้ผู้ใช้ 41 คนประเมินโดยไม่รู้ชื่อโมเดล และนำตัวเลือกที่ดีที่สุดไปพัฒนาเป็นเว็บที่มีการทดลองใช้จริง 100 ครั้ง

\subsection{สรุปข้อค้นพบหลัก}

ข้อค้นพบแรกคือการปรับจูนได้ผลกับทั้งสองโมเดล แต่ได้ผลไม่เท่ากัน
HairFastGAN Trial B Epoch 24 (B24) ลดความคลาดเคลื่อนของสีผมในปริภูมิสี CIELAB (Lab) จาก 13.0 เหลือ 9.0
เพิ่มความเหมือนของใบหน้าที่วัดด้วย ArcFace (Additive Angular Margin Loss) จาก 0.73 เป็น 0.79
ลดค่าความคลาดเคลื่อนสัมบูรณ์เฉลี่ย (Mean Absolute Error: MAE) นอกบริเวณผมจาก 0.041 เหลือ 0.014
และเพิ่มดัชนีความคล้ายเชิงโครงสร้าง (Structural Similarity Index: SSIM) จาก 0.904 เป็น 0.947
ผลนี้แก้ปัญหาสีเป็นปื้นและใบหน้าซีดที่พบตั้งแต่ต้นได้ตรงจุด
ส่วน Barbershop-FT (Barbershop ที่ปรับจูนแล้ว) Trial A Epoch 30 (A30) ดีขึ้นน้อยกว่า แต่ดีขึ้นอย่างสม่ำเสมอ
เช่น ArcFace เพิ่มขึ้น $+0.013$ (ช่วงความเชื่อมั่น 95\% คือ $+0.008$ ถึง $+0.019$) และดีขึ้นครบ 10 จาก 10 ภาพ
ผู้วิจัยตั้งข้อสังเกตว่าความต่างนี้อาจเกี่ยวข้องกับตำแหน่งของส่วนที่ปรับได้
แต่ยังไม่ได้แยกพิสูจน์ด้วยการทดลองตัดองค์ประกอบออกทีละส่วน (Ablation Study)
HairFastGAN มีโมดูลที่ตรงกับปัญหาให้ปรับ ขณะที่ Barbershop ปรับได้เฉพาะชั้นวาดภาพ ซึ่งอยู่หลังขั้นที่กำหนดทรงไปแล้ว

ข้อค้นพบที่สองคือแต่ละโมเดลเด่นกันคนละด้าน
HairFastGAN รักษาใบหน้าและบริเวณที่ไม่ควรเปลี่ยนได้ดีที่สุด และเร็วกว่าประมาณ 65 เท่า
(ประมาณ 10 วินาที เทียบกับ 657 วินาทีต่อภาพ)
ส่วน Barbershop-FT ให้สีผมใกล้เป้าหมายที่สุด (ความคลาดเคลื่อน 7.86) แต่ทรงผมแทบไม่เปลี่ยนจากก่อนปรับจูน

ข้อค้นพบที่สามคือกติกาการเลือกที่มีด่านกันความผิดพลาดมีความจำเป็นจริง
Barbershop-FT Trial C และ D มีคะแนนรวมใกล้เคียงกับ Trial A
แต่ค่าการทับซ้อนของทรงผม (Intersection over Union: IoU) ต่ำกว่าตอนก่อนฝึกทุก Epoch
ถ้าใช้คะแนนอย่างเดียว อาจได้โมเดลที่ทรงผมแย่ลงโดยไม่รู้ตัว

ข้อค้นพบที่สี่คือ Diffusion แบบแรงต่ำทำให้ภาพดูเป็นธรรมชาติขึ้น แต่มีสิ่งที่ต้องแลก
ผู้ใช้เห็นว่าภาพของ Barbershop-FT สมจริงขึ้น $+0.59$ (95\% CI $+0.27$ ถึง $+0.93$)
คะแนนเฉลี่ยความพึงพอใจเพิ่มขึ้น $+0.20$ ถึง $+0.24$ และคะแนนเฉลี่ยของทรงกับสีผมไม่ลดลง
แต่ช่วงความเชื่อมั่นของความพึงพอใจยังแตะหรือครอบคลุมเลข 0
สิ่งที่ต้องแลกคือคะแนนเฉลี่ยความเหมือนของใบหน้าที่ลดลง $-0.22$ ถึง $-0.32$ (ช่วงความเชื่อมั่นยังแตะหรือครอบคลุมเลข 0)
ซึ่งไปทางเดียวกับ ArcFace ที่ลดลง 0.03--0.19 ทุกกรณีในกรณีศึกษา 8 ชุดภาพของ HFG B24 และ BS-FT A30
ขณะที่พื้นที่นอกหน้ากาก (Mask) เปลี่ยนน้อยในเกือบทุกชุด (MAE เฉลี่ย 0.0045 สำหรับ HFG และ 0.0065 สำหรับ BS-FT)

ข้อค้นพบที่ห้า ซึ่งผู้วิจัยคิดว่าสำคัญที่สุด คือตัวเลขอย่างเดียวตัดสินคุณภาพของภาพที่สร้างขึ้นใหม่ไม่ได้
ในกรณีศึกษา 8 ชุดภาพ หลังผ่าน Diffusion คะแนนคุณภาพภาพ MUSIQ (Multi-scale Image Quality Transformer) ไม่ดีขึ้นถึงเกณฑ์ และ ArcFace ลดลง
แต่ผู้ใช้กลับบอกว่าภาพดีขึ้น
MUSIQ ประเมินคุณภาพจากภาพผลลัพธ์โดยไม่ใช้ภาพอ้างอิง ส่วน ArcFace เปรียบเทียบเวกเตอร์ใบหน้ากับภาพต้นฉบับ
ข้อมูลในกรณีศึกษานี้ยังไม่พอจะบอกว่าทำไมตัวเลขกับคะแนนผู้ใช้จึงไปคนละทาง และไม่ได้แปลว่าตัวชี้วัดเหล่านี้ใช้ไม่ได้
ถ้าดูเพียงตัวเลข ผู้วิจัยคงตัด Diffusion ออกไปแล้ว
จึงควรใช้ตัวชี้วัดและการประเมินโดยผู้ใช้ร่วมกัน และตีความตามขอบเขตของกลุ่มตัวอย่างที่วัด

ข้อค้นพบที่หกคือความเร็วเป็นตัวตัดสินในการใช้งานจริง
ร้อยละ 84 ของคำตอบยอมรอได้ไม่เกิน 5 นาที แม้ Barbershop-FT + Diffusion จะสมจริงที่สุด
แต่ใช้เวลาประมาณ 11 นาที ผู้วิจัยจึงเลือก HairFastGAN B24 + Diffusion
ซึ่งได้คะแนนความพึงพอใจในการใช้งานจริงสูงสุด (3.90) สีผมตรงที่สุด (4.22) ใช้เวลาประมาณ 13 วินาที
และรักษาใบหน้าได้ดีกว่า Barbershop-FT + Diffusion $+0.73$ (95\% CI $+0.37$ ถึง $+1.10$)

ข้อค้นพบสุดท้ายคือผู้ใช้พอใจเมื่อได้ใช้งานจริง การทดลองใช้เว็บ 100 ครั้งได้คะแนนความพึงพอใจโดยรวม 4.73 จาก 5
(ร้อยละ 92 ให้ 4 หรือ 5 คะแนน) คะแนนความพึงพอใจต่อเวลารอ 4.88
และเวลารอจริงมีค่ามัธยฐาน 14.9 วินาที โดยร้อยละ 85 เสร็จภายใน 20 วินาที

\subsection{การบรรลุวัตถุประสงค์ของโครงงาน}

โครงงานบรรลุวัตถุประสงค์ทั้ง 3 ข้อ
วัตถุประสงค์ข้อแรก คือการศึกษาและเปรียบเทียบโมเดล ทำได้โดยเปรียบเทียบ HairFastGAN และ Barbershop
ทั้งก่อนและหลังปรับจูนภายใต้ข้อมูลและตัวชี้วัดชุดเดียวกัน
พร้อมระบุให้ชัดถึงค่าสูญเสียส่วนเพิ่มของ HairFastGAN ความต่างของการบันทึกน้ำหนัก และการแก้ไขเกณฑ์การเลือก
วัตถุประสงค์ข้อที่สอง คือการวิเคราะห์ข้อดีและข้อจำกัด ทำได้โดยวัดด้วยตัวชี้วัด 5 ด้าน เวลาประมวลผล
ผลของขั้น Diffusion และความเห็นของผู้ใช้ 41 คน ซึ่งทำให้เห็นทั้งจุดแข็ง จุดอ่อน และสิ่งที่ต้องแลกของแต่ละตัวเลือก
และวัตถุประสงค์ข้อที่สาม คือการพัฒนาเว็บแอปพลิเคชันต้นแบบ ทำได้โดยใช้ HairFastGAN B24 + Diffusion
และมีการทดลองใช้งานจริง 100 ครั้ง ได้คะแนนความพึงพอใจโดยรวม 4.73 จาก 5

% -----------------------------------------------------------
\section{ข้อเสนอแนะ}
% -----------------------------------------------------------

ถ้าต้องการให้ทรงผมของ Barbershop ดีขึ้น ควรเพิ่มจำนวนรอบในการค้นหา โดยเฉพาะขั้น Alignment ที่กำหนดทรงผม (ปัจจุบัน 140 รอบ)
แต่เวลาประมวลผลซึ่งนานอยู่แล้ว (ประมาณ 657 วินาทีต่อภาพ) จะยิ่งเพิ่มขึ้นอีก

เรื่องใบหน้าที่เปลี่ยนหลังผ่าน Diffusion ควรลดความแรงของ Diffusion เฉพาะบริเวณใบหน้า
หรือให้ Diffusion แก้ไขเฉพาะบริเวณผม เพื่อให้ใบหน้ายังเหมือนเดิม แต่ผมยังชัดขึ้น

เรื่องข้อมูล ชุดตรวจสอบและชุดทดสอบมีจำนวนน้อย (9 และ 10 ชุด) ทำให้ผลบางด้าน เช่น สีผมของ HairFastGAN ยังแกว่ง
ควรเพิ่มจำนวนชุดทดสอบ และสำหรับ Barbershop-FT ควรบันทึกน้ำหนักไว้ทุก Epoch
เพื่อให้เลือก Checkpoint ได้จากทุก Epoch ไม่ใช่เพียง Epoch 30

เรื่องเว็บ LOOKLAB ผู้ใช้บอกว่าผมบางทรงฟูเกินจริงและไม่ค่อยเรียบ ควรปรับปรุงส่วนที่กำหนดทรงผม
และควรลดเวลารอในโหมดกล้อง ซึ่งตอนนี้ยังต้องรอประมาณ 15 วินาที ให้ใกล้เคียงแบบเรียลไทม์มากขึ้น
